\section{Materials and Methods}\label{sec:methods}

\begin{figure*}[!t]
    \centering
    \includegraphics[width=\textwidth]{figures/piHurdleSpec_editable (4).pdf}
    \caption{
        Overview of the two-stage \modelname{} framework for top-down tandem
        mass spectrum prediction. (a) Overall sequence-to-spectrum pipeline.
        (b) Stage 1 learns fragment-ion occurrence and positive-intensity
        distributions. (c) Stage 2 maps the learned distribution parameters
        to a point spectrum using an independent distribution-to-spectrum
        decoder. (d) Charge- and cleavage-resolved output representation.
    }
    \label{fig:hurdlespec_framework}
\end{figure*}

\subsection{Overview of the \modelname{} Architecture}
\label{sec:architecture_overview}

Top-down tandem mass spectrum prediction maps a proteoform sequence and its
experimental conditions to the occurrence and relative intensity of candidate
fragment ions. The output is highly sparse along both the cleavage-site and
charge-state dimensions: most theoretically valid ions are not observed, and
the intensities of observed ions vary substantially. Direct regression over
all entries therefore mixes two distinct events, namely whether an ion is
observed and how intense it is when observed.

\modelname{} addresses this problem with a hurdle-modeling strategy that
separates fragment-ion occurrence from positive fragment-ion intensity
(Fig.~\ref{fig:hurdlespec_framework}). The model operates on a structured
charge-resolved spectrum rather than an unstructured list of peaks. For a
proteoform of length $L$, the target is
\begin{equation}
    \mathbf{Y}\in\mathbb{R}_{\geq 0}^{(L-1)\times K\times Z},
\end{equation}
where $K=2$ denotes the N-terminal and C-terminal ion series and $Z=30$ is the
maximum fragment charge. Each entry $y_{t,k,z}$ represents the relative
intensity at cleavage $t$, ion series $k$, and charge $z$.

Training comprises two stages. Stage 1 learns an explicit distributional
representation $(\boldsymbol{\pi},\boldsymbol{\mu},\mathbf{s},\mathbf{U})$.
Here, $\boldsymbol{\pi}$ contains occurrence probabilities,
$\boldsymbol{\mu}$ contains conditional positive-intensity means, and
$(\mathbf{s},\mathbf{U})$ parameterizes the conditional covariance matrix.
Stage 2 introduces an architecturally independent distribution-to-spectrum
decoder. The decoder receives only these distribution parameters, without the
amino-acid sequence or encoder hidden states, and directly predicts the final
charge-resolved spectrum. During Stage 2, the encoder, distribution heads, and
decoder are jointly fine-tuned using a spectral-angle objective together with
a weak distribution-preserving regularizer.

\subsection{Data Preprocessing}
\label{sec:data_preprocessing}

We used the annotated training and validation partitions of the TDRepo
top-down tandem mass spectrometry data. Proteoforms longer than 200 residues,
precursor charge states above 30, peptide-spectrum matches with an E-value
greater than $10^{-4}$, and spectra with charge-resolved fragment coverage
below 0.2 were excluded. The resulting training and validation sets contained
601,520 and 152,124 spectra, respectively. Proteoform sequences did not overlap
between the two partitions.

Each spectrum was converted to a tensor of shape $199\times2\times30$.
N-terminal $b/c$ ions and C-terminal $y/z$ ions were assigned to the two ion
series. An ion $b_i$ or $c_i$ was mapped to cleavage index $i-1$, whereas an
ion $y_i$ or $z_i$ was mapped to index $L-i-1$. Fragment annotations whose
charge exceeded the precursor charge were discarded. Isotopic annotations
with offsets $-1$, $0$, and $+1$ were folded into the corresponding canonical
$(t,k,z)$ entry following the TDRepo preprocessing procedure.

Fragment intensities were divided by the maximum annotated charge-resolved
fragment intensity within each spectrum. No logarithmic transformation was
applied. We additionally constructed an occurrence tensor
$\mathbf{O}=\mathbb{I}(\mathbf{Y}>0)$ and a theoretical-ion mask
$\mathbf{M}$, where $m_{t,k,z}=1$ only when the cleavage and charge combination
is valid for the corresponding proteoform. Padding positions and invalid
charge states were excluded from all objectives and evaluation metrics.

\subsection{Distribution Learning}
\label{sec:distribution_learning}

\subsubsection{Sequence and Experimental Metadata Encoding}
\label{sec:sequence_encoding}

Let $\mathbf{a}=(a_1,\ldots,a_L)$ denote a proteoform sequence and let cleavage
$t$ lie between residues $a_t$ and $a_{t+1}$. We represent each cleavage by
the two adjacent residues together with spectrum-level experimental metadata.
The metadata comprise the fragmentation method, instrument type, collision
energy, and precursor charge. Collision energy is discretized into ten
intervals of width 10, with an additional category for a missing value. The
architecture also reserves an embedding for the N-terminal form; this field
is unavailable in the current TDRepo preprocessing and is therefore assigned
the same unknown category to all spectra.

Each amino acid is mapped to a 64-dimensional embedding, and each metadata
field is mapped to a 16-dimensional embedding. Spectrum-level metadata
embeddings are broadcast to all valid cleavage positions. The input vector at
cleavage $t$ is
\begin{equation}
\begin{split}
    \mathbf{x}_t = [&\,\mathbf{e}_{\mathrm{aa}}(a_t)
    \mathbin\Vert \mathbf{e}_{\mathrm{aa}}(a_{t+1})
    \mathbin\Vert \mathbf{e}_{\mathrm{frag}}(f)
    \mathbin\Vert \mathbf{e}_{\mathrm{inst}}(q)\\
    &\mathbin\Vert \mathbf{e}_{\mathrm{energy}}(e)
    \mathbin\Vert \mathbf{e}_{\mathrm{prec}}(z_p)
    \mathbin\Vert \mathbf{e}_{\mathrm{Nterm}}(n)],
\end{split}
\label{eq:cleavage_input}
\end{equation}
where $\mathbin\Vert$ denotes concatenation. Thus,
$\mathbf{x}_t\in\mathbb{R}^{208}$, consisting of two 64-dimensional residue
embeddings and five 16-dimensional metadata embeddings.

The cleavage sequence is layer-normalized and processed by a residual local
convolutional module with successive kernel sizes 3 and 5, yielding an
effective receptive field of seven cleavage positions. This module captures
local sequence contexts associated with fragmentation. Its output is passed
to a two-layer bidirectional LSTM with 256 hidden units in each direction to
model longer-range dependencies along the proteoform. The final encoder state
$\mathbf{h}_t\in\mathbb{R}^{512}$ is used by both the occurrence and
positive-intensity distribution heads. Padding positions are excluded by a
cleavage mask, and only valid sequence lengths are processed by the recurrent
encoder.

\subsubsection{Fragment-Ion Occurrence Modeling}
\label{sec:occurrence_modeling}

The occurrence head maps each $\mathbf{h}_t$ through a 128-dimensional hidden
layer and predicts 60 logits, corresponding to two ion series and 30 fragment
charges. The occurrence probability for candidate ion $j$ is
\begin{equation}
    \pi_j = \sigma(a_j),
\end{equation}
where $a_j$ is the predicted logit. We optimize the binary cross-entropy over
all theoretically valid ions:
\begin{equation}
    \mathcal{L}_{\mathrm{BCE}}
    = -\frac{1}{\sum_j m_j}
    \sum_j m_j
    \left[o_j\log\pi_j+(1-o_j)\log(1-\pi_j)\right].
\end{equation}
Thus, both observed and unobserved theoretical ions contribute to occurrence
learning.

\subsubsection{Fragment-Ion Intensity Distribution Modeling}
\label{sec:intensity_distribution}

The positive-intensity head predicts a conditional mean $\boldsymbol{\mu}$
and a low-rank covariance parameterization. For each spectrum, the conditional
covariance over charge-resolved fragment-ion variables is
\begin{equation}
    \boldsymbol{\Sigma}
    = \operatorname{diag}(\mathbf{s}^{2})
      + \mathbf{U}\mathbf{U}^{\top},
    \label{eq:low_rank_covariance}
\end{equation}
where $\mathbf{s}$ contains ion-specific residual scales and
$\mathbf{U}\in\mathbb{R}^{D\times r}$ captures shared variation among ions.
We use rank $r=4$. The mean and scale are bounded as
\begin{align}
    \mu_j &= 10^{-5} + (1-10^{-5})\sigma(b_j),\\
    s_j &= 10^{-3} + (0.5-10^{-3})\sigma(c_j),
\end{align}
and each low-rank factor is constrained by
$U_{jq}=0.1\tanh(d_{jq})$.

Let $\mathcal{O}_i=\{j:y_{ij}>0,\,m_{ij}=1\}$ be the observed positive-ion set
of spectrum $i$. We restrict the predicted mean and covariance to this set and
minimize
\begin{equation}
\begin{split}
    \mathcal{L}_{\mathrm{MVN}}
    = \frac{1}{\sum_i |\mathcal{O}_i|}
      \sum_i \frac{1}{2}\Big[
      &|\mathcal{O}_i|\log(2\pi)
      +\log|\boldsymbol{\Sigma}_{i,\mathcal{O}_i}|\\
      &+(\mathbf{y}_{i,\mathcal{O}_i}-\boldsymbol{\mu}_{i,\mathcal{O}_i})^{\top}
      \boldsymbol{\Sigma}_{i,\mathcal{O}_i}^{-1}
      (\mathbf{y}_{i,\mathcal{O}_i}-\boldsymbol{\mu}_{i,\mathcal{O}_i})
      \Big].
\end{split}
\label{eq:mvn_loss}
\end{equation}
Zeros are excluded from this likelihood and are modeled only by the Bernoulli
component. Equation~\ref{eq:mvn_loss} is a Gaussian working likelihood for
observed positive intensities; it does not assume that a single fixed
parametric marginal distribution is shared by all ion types and conditions.
The matrix determinant lemma and Woodbury identity reduce the required matrix
operations to the $r\times r$ low-rank system.

The Stage-1 objective is
\begin{equation}
    \mathcal{L}_{\mathrm{dist}}
    = 0.1\mathcal{L}_{\mathrm{MVN}}
      +0.5\mathcal{L}_{\mathrm{BCE}}.
    \label{eq:stage1_loss}
\end{equation}

\subsubsection{Expectation-Based Statistical Prediction}
\label{sec:expectation_prediction}

The hurdle representation provides a direct statistical point estimate
\begin{equation}
    \widetilde{y}_j
    = \mathbb{E}[Y_j\mid\mathbf{x}]
    = \pi_j\mu_j.
    \label{eq:hurdle_expectation}
\end{equation}
We use this expectation as the Stage-1 prediction for diagnostic evaluation.
It is not supplied to the Stage-2 decoder as an input or used as a residual
prediction baseline; the decoder instead receives
$(\boldsymbol{\pi},\boldsymbol{\mu},\mathbf{s},\mathbf{U})$ separately.

\subsection{Task Adaptation}
\label{sec:task_adaptation}

\subsubsection{Distribution-to-Spectrum Decoder}
\label{sec:spectrum_decoder}

The Stage-2 decoder converts the complete distributional representation into
a point spectrum. At each cleavage, the 60 occurrence probabilities, 60
conditional means, 60 residual scales, and $60\times4$ low-rank factors are
concatenated into a 420-dimensional feature vector. Therefore, the decoder can
use occurrence, location, marginal uncertainty, and cross-ion covariance
information while remaining isolated from the original sequence features and
encoder states.

The 420-dimensional vector is layer-normalized and projected to 192 hidden
units. A masked mean over cleavage positions is processed by a two-layer MLP
and added to every position to provide spectrum-level context. Three temporal
residual blocks then exchange information among neighboring cleavage sites.
Each block consists of layer normalization, a width-3 convolution from 192 to
384 channels, a gated linear unit, GELU activation, dropout, and a width-1
projection back to 192 channels. A final linear head predicts 60 intensity
logits at each cleavage. The output spectrum is
\begin{equation}
    \widehat{\mathbf{Y}}
    = \operatorname{MaxNorm}\!\left(
        \operatorname{Softplus}\!\left[
        f_{\theta}(\boldsymbol{\pi},\boldsymbol{\mu},\mathbf{s},\mathbf{U})
        \right]\odot\mathbf{M}
      \right),
    \label{eq:decoder_prediction}
\end{equation}
where $\operatorname{MaxNorm}$ divides all charge-resolved intensities in a
spectrum by their maximum. This direct mapping distinguishes the decoder from
a linear calibration or multiplicative correction of
$\boldsymbol{\pi}\odot\boldsymbol{\mu}$.

\subsubsection{Spectral Similarity Objective}
\label{sec:spectral_similarity}

The decoder is optimized with a spectral-angle loss over all theoretically
valid charge-resolved entries. After flattening the masked target and
prediction, their cosine similarity is
\begin{equation}
    c_i =
    \frac{\langle\widehat{\mathbf{y}}_i,\mathbf{y}_i\rangle}
    {\|\widehat{\mathbf{y}}_i\|_2\|\mathbf{y}_i\|_2}.
\end{equation}
The loss is
\begin{equation}
    \mathcal{L}_{\mathrm{SA}}
    = \frac{1}{N}\sum_{i=1}^{N}
      \frac{2}{\pi}\arccos(c_i).
    \label{eq:sa_loss}
\end{equation}
Unlike a pointwise regression loss, this objective directly optimizes the
relative spectral pattern used for spectrum comparison.

During Stage 2, the distribution objective is retained with a small weight:
\begin{equation}
    \mathcal{L}_{\mathrm{stage2}}
    = \mathcal{L}_{\mathrm{SA}}
      +0.01\mathcal{L}_{\mathrm{dist}}.
    \label{eq:stage2_loss}
\end{equation}
The decoder, distribution heads, and sequence encoder are all trainable in
this stage. The retained term discourages the distribution parameters from
degenerating into unconstrained intermediate features under the task loss.

\subsection{Optimization and Training Strategy}
\label{sec:optimization}

Stage 1 was trained for at most 20 epochs with early stopping of five epochs
based on validation distribution loss. The sequence backbone and distribution
heads used learning rates of $10^{-4}$ and $5\times10^{-4}$, respectively.
The selected Stage-1 checkpoint was obtained at epoch 12. Stage 2 was trained
for 30 epochs with a batch size of 2048. We used AdamW with weight decay
$10^{-4}$ and separate learning rates of $5\times10^{-5}$,
$10^{-4}$, and $5\times10^{-4}$ for the sequence backbone, distribution heads,
and decoder, respectively. A cosine annealing schedule reduced the learning
rates to $10^{-5}$. Gradients were clipped to a norm of 5, and mixed-precision
training was used. The Stage-2 checkpoint with the lowest validation
$\mathcal{L}_{\mathrm{SA}}$ was selected. All experiments used random seed
1337.

The current implementation initializes the Stage-1 sequence backbone from a
previously trained NovaPredict intensity-prediction checkpoint; newly added
distribution heads are initialized independently. This initialization must be
applied consistently or reported explicitly when comparing against models
trained from scratch.

\subsection{Evaluation Metrics}
\label{sec:evaluation_metrics}

We report two complementary cosine-similarity metrics. Charge-resolved cosine
similarity compares the flattened tensors over cleavage site, ion series, and
fragment charge:
\begin{equation}
    \operatorname{Cos}_{\mathrm{resolved}}(\widehat{\mathbf{Y}},\mathbf{Y})
    = \cos\!\left(\operatorname{vec}(\widehat{\mathbf{Y}}\odot\mathbf{M}),
                  \operatorname{vec}(\mathbf{Y}\odot\mathbf{M})\right).
\end{equation}
Charge-summed cosine similarity first sums intensities over fragment charges
and then compares the resulting cleavage-by-ion-series matrices:
\begin{equation}
    \operatorname{Cos}_{\mathrm{summed}}(\widehat{\mathbf{Y}},\mathbf{Y})
    = \cos\!\left(
        \operatorname{vec}\!\left(\sum_z\widehat{\mathbf{Y}}_{:,:,z}\right),
        \operatorname{vec}\!\left(\sum_z\mathbf{Y}_{:,:,z}\right)
      \right).
\end{equation}
Both metrics are computed per spectrum and then averaged across the evaluation
set. The first evaluates charge-state assignment and intensity jointly,
whereas the second evaluates the predicted fragmentation pattern after
collapsing charge states.
